\section{Discussão}
\label{sec:discussao}

\subsection{As hipóteses da Introdução contra os resultados}

\begin{table}[ht]
\centering
\caption{As oito hipóteses informais registradas antes da análise, uma por RQ.}
\label{tab:hipoteses}
\small
\begin{tabular}{p{0.7cm}p{4.4cm}p{5.4cm}l}
\toprule
 & \textbf{Hipótese} & \textbf{Resultado} & \textbf{Veredito} \\
\midrule
H1 & Frequência alta, em \emph{Elite}/\emph{High} &
     Mediana de 0,46 release/semana (\emph{Medium}); 6 \emph{Elite} e 69
     \emph{High} contra 152 \emph{Medium} e 73 \emph{Low} & Refutada \\
H2 & As variantes de lead time são próximas &
     Medianas de 324,9 h e 411,8 h; $\rho = 0{,}71$, mas IQR de 880 h contra 1.716 h &
     Parcial \\
H3 & CFR baixo e variantes concordantes &
     CFR (a) mediana 0,056 (\emph{Elite}), mas as variantes discordam em 72,3\% dos
     repositórios, com $\kappa = -0{,}015$ & Parcial \\
H4 & Recuperação abaixo de um dia &
     Mediana de 4,94 h; 181 dos 300 em \emph{High} e 33 em \emph{Elite} &
     Confirmada \\
H5 & Sem \emph{trade-off} velocidade × estabilidade &
     Confirma-se com o CFR (a) ($\rho = 0{,}07$); com o CFR (b), $\rho = 0{,}62$ por
     artefato de construção & Parcial \\
H6 & Estrelas associadas a melhor desempenho &
     Nenhuma das quatro métricas se associa a estrelas após Holm; linguagem e
     contribuidores, sim & Refutada \\
H7 & Conclusões robustas à definição &
     Até 31,3\% mudam de categoria; trocar a variante pesa quase 4× mais que trocar a
     unidade & Refutada \\
H8 & Retrabalho baixo e próximo do CFR (b) &
     Mediana de 0,70 (alta) e correlação apenas fraca com o CFR (b)
     ($\rho = 0{,}36$) & Refutada \\
\bottomrule
\end{tabular}
\end{table}

\subsection{Velocidade e Lead Time (H1 e H2)}

A H1 foi \textbf{refutada}. A mediana de 0,46 release por semana equivale a uma entrega a
cada 15 dias e enquadra a amostra em \emph{Medium}, não em \emph{Elite} ou \emph{High}:
apenas 6 dos 300 repositórios publicam diariamente ou mais. Projetos open-source populares
operam em regime \emph{descontínuo} de releases, agrupando mudanças em lotes, mesmo quando
dispõem de automação madura de CI.

A H2 foi \textbf{parcialmente confirmada}, e o modo como ela falha é informativo. As duas
variantes correlacionam-se fortemente ($\rho = 0{,}71$), o que sustenta a expectativa de que
medem o mesmo fenômeno; mas divergem em escala e, sobretudo, em dispersão: a variante (b)
tem IQR de 1.716 h contra 880 h da (a). A razão é estrutural. A variante (a) é definida pelo
commit \emph{mais antigo} de cada release, enquanto a (b) é a mediana de \emph{todos} os
commits. Basta um commit antigo esquecido no intervalo --- um \emph{branch} retomado, um
\emph{cherry-pick}, um \emph{pull request} parado --- para a (a) subir sem que o ritmo de
trabalho tenha mudado. As duas respondem a perguntas distintas: a (a), "quanto tempo cabe
dentro de uma entrega"; a (b), "quanto tempo um commit típico espera para ser entregue".

\subsection{A Ilusão de Estabilidade do \emph{Proxy} de CI (H3)}

A H3 foi \textbf{parcialmente confirmada}: a taxa é de fato baixa pelo \emph{proxy} de CI
(mediana de 5,63\%, faixa \emph{Elite}, com 234 dos 300 repositórios nessa faixa), mas a
previsão de que as duas variantes concordariam está errada por larga margem --- elas
discordam de faixa em 72,3\% dos repositórios, com $\kappa = -0{,}015$.

A diferença não é de calibração, é de objeto. Uma execução de CI vermelha é, na maior parte
das vezes, um teste quebrado em um \emph{pull request}: barata, detectada antes da entrega e
frequentemente a prova de que o CI está fazendo o seu trabalho. Uma release seguida de uma
corretiva em até sete dias é um defeito que atravessou o pipeline e chegou ao usuário. O
\emph{proxy} de CI mede a \emph{barreira}; o de entrega mede o que \emph{passou} por ela.
Comparar qualquer um dos dois diretamente com as faixas do relatório DORA, que falam de
produção, induz a erro --- e o CFR (a), com 78\% da amostra em \emph{Elite}, sugere uma
estabilidade que não mediu.

\subsection{O Artefato de Construção do CFR de Entrega (H4 e H5)}

A H4 foi \textbf{confirmada}: a recuperação mediana de 4,94 h fica bem abaixo das 24 h
previstas, com 181 dos 300 repositórios em \emph{High} e 33 em \emph{Elite}. Uma \emph{build}
quebrada no \emph{default branch} é, de fato, tratada com prioridade.

A H5 é o ponto mais delicado do estudo. Com o \emph{proxy} de CI, a ausência de
\emph{trade-off} se confirma ($\rho = 0{,}07$, não significativo). Com o \emph{proxy} de
entrega, a correlação é positiva e moderada ($\rho = 0{,}62$), o que, lido ingenuamente,
contradiria o DORA e sugeriria que entregar mais causa mais falhas. O diagnóstico da janela
de sete dias mostra que a leitura é um \textbf{artefato de construção da própria métrica}:
no quartil superior de frequência, 87,7\% dos intervalos entre releases cabem em sete dias,
contra 25,0\% no quartil inferior. Quem publica a cada 1,4 dia quase sempre tem uma release
seguinte dentro da janela de observação da heurística e, portanto, muito mais chance de ser
marcado como falha; quem publica a cada 28 dias quase nunca tem. O numerador e o denominador
compartilham a mesma causa --- a cadência --- e o coeficiente mede cadência disfarçada de
estabilidade. Correlação não implica causalidade; aqui, sequer há correlação legítima a
interpretar.

\subsection{Fatores Estruturais e o Peso da Definição (H6, H7 e H8)}

A H6 foi \textbf{refutada} no ponto que previa: estrelas não se associam a nenhuma das
quatro métricas após a correção de Holm. A popularidade parece operar como indicador de
reputação, desacoplado do fluxo interno de engenharia. O que sobrevive são linguagem
($\varepsilon^2 = 0{,}164$ no lead time) e número de contribuidores
($\varepsilon^2 = 0{,}100$), ambos com efeito médio. Vale separar significância de
relevância: com $n = 300$, efeitos pequenos alcançam significância sem significar muito.

A H7 foi \textbf{refutada}, e o desenho das combinações permite dizer \emph{como}. Entre C2
e C4, que diferem apenas na unidade de deploy, somente 7,3\% dos repositórios mudam de
categoria ($\kappa = 0{,}878$). Entre C1 e C4, que diferem apenas nas variantes das métricas,
28,3\% mudam ($\kappa = 0{,}546$). Ou seja: \textbf{trocar a variante da métrica desloca a
classificação quase quatro vezes mais do que trocar a unidade de deploy}. A discussão sobre
se \emph{tag} ou release conta como deploy, que ocupa boa parte da literatura de mineração,
importa menos do que a escolha entre CFR de CI e CFR de entrega.

A H8 foi \textbf{refutada} nas duas partes: a taxa de retrabalho não é baixa (mediana de
70,0\%) e sua correlação com o CFR de entrega é apenas fraca ($\rho = 0{,}36$), embora ambas
derivem da mesma heurística. A explicação está na janela: o retrabalho conta toda release
corretiva do período, enquanto o CFR (b) só conta aquelas que \emph{seguem} outra release em
até sete dias. Repositórios que corrigem com calma produzem retrabalho alto e CFR (b) baixo.

\subsection{A Fragilidade das Fronteiras}

Soma-se a tudo isso um achado que não estava previsto em nenhuma hipótese: 86,3\% dos
repositórios mudam de categoria geral se uma única métrica avançar um degrau. A
classificação em faixas é uma discretização de medidas contínuas, e nada no método sinaliza
quem está na fronteira. Esse resultado nasceu de uma divergência relatada na replicação
cruzada (Seção \ref{sec:replicacao}) e reposiciona a interpretação de qualquer comparação de
categorias entre estudos.
